\chapter{Appendices}\label{sec:appendices}
\section{Appendix A: Mathematical Notation}\label{sec:appendix-a-mathematical-notation}

{\def\LTcaptype{table} % do not increment counter
\begin{longtable}[]{@{}ll@{}}
\toprule\noalign{}
Symbol & Description \\
\midrule\noalign{}
\endhead
\bottomrule\noalign{}
\endlastfoot
\(b\) & Binary bit (\(\in \{0, 1\}\)) \\
\(x\) & Modulated BPSK symbol (\(\in \{-1, +1\}\)) \\
\(y\) & Received signal \\
\(n\) & Noise sample \\
\(h\) & Fading coefficient \\
\(\sigma^2\) & Noise variance \\
\(\text{SNR}\) & Signal-to-Noise Ratio (linear) \\
\(G\) & AF amplification gain \\
\(\mathbf{W}\) & Neural network weight matrix \\
\(\mathbf{b}\) & Bias vector \\
\(P_e\) & Bit error rate \\
\(Q(\cdot)\) & Gaussian Q-function \\
\(\beta\) & VAE KL weight \\
\(\lambda\) & Regularization parameter \\
\(\eta\) & Learning rate \\
\(\Delta\) & State space discretization step \\
\(\mathbf{A}, \mathbf{B}, \mathbf{C}, D\) & State space model matrices \\
\end{longtable}
}

\section{Appendix B: Model Architectures and Hyperparameters}\label{sec:appendix-b-model-architectures-and-hyperparameters}

% column helpers
\newcolumntype{L}[1]{>{\raggedright\arraybackslash}p{#1}}
\newcolumntype{Y}{>{\raggedright\arraybackslash}X}

\begin{table}[ht]
\centering
\footnotesize
\setlength{\tabcolsep}{4pt}
\renewcommand{\arraystretch}{1.15}

\begin{tabularx}{\textwidth}{|L{2.1cm}|Y|L{3.2cm}|L{1.7cm}|L{2.3cm}|}
\hline
\textbf{Relay} &
\textbf{Network Topology} &
\textbf{Train Loss (LR)} &
\textbf{Total Params} &
\textbf{Implementation} \\
\hline
MLP (Minimal) &
Feedforward MLP, 5-symbol sliding window &
MSE, lr=0.01, 100 epochs &
169 &
NumPy (CPU) \\
\hline
Hybrid &
MLP + DF switching via learned SNR threshold ($\sim$5 dB) &
MSE (MLP only), lr=0.01 &
169 &
NumPy (CPU) \\
\hline
VAE &
Encoder--decoder latent model (7$\to$8$\to$1) &
$\beta$-VAE ($\beta=0.1$), Adam lr=$10^{-3}$ &
1{,}777 &
PyTorch (CUDA) \\
\hline
cGAN (WGAN-GP) &
Conditional GAN (generator + critic) &
WGAN-GP, Adam lr=$10^{-4}$ &
2{,}946 &
PyTorch (CUDA) \\
\hline
Transformer &
2-layer encoder, 4 heads, $d_{\text{model}}=32$ &
MSE, Adam lr=$10^{-3}$ &
17{,}697 &
PyTorch (CUDA/CPU) \\
\hline
Mamba-S6 &
Selective SSM, 2 S6 layers &
MSE, Adam lr=$10^{-3}$ &
24{,}001 &
PyTorch (CUDA/CPU) \\
\hline
Mamba-2 SSD &
Selective SSM with SSD kernel (2 layers) &
MSE, Adam lr=$10^{-3}$, clip=1.0 &
26{,}179 &
PyTorch (CUDA) \\
\hline
\end{tabularx}

\caption[Relay models, objectives and parameter counts]{Comparison of relay models, training objectives, parameter counts, and implementations.}
\label{tab:relay_comparison}
\end{table}

\section{Appendix C: Reproducibility}\label{sec:appendix-c-software-architecture}

\subsection{Testing} 159 automated tests (pytest) cover the channels, modulation, relay-strategy, simulation and statistics modules with a 100\% pass rate. The suite requires only NumPy and SciPy; PyTorch is optional and needed solely by the neural relays that use it.

\subsection{Reproducibility} Random seeds are controlled at the source (bit generation) and noise (per-trial seeding) levels to ensure reproducible results.

\paragraph{Provenance of the reported results.} Every result reported in this thesis was regenerated from the archived experiment artifacts and verified against its committed result file, at $10$ trials $\times\ 10{,}000$ bits per SNR point. The verification is mechanical rather than manual: each numerical table cell is re-derived from the committed result file it came from, so a table and the data behind it cannot drift apart unnoticed.



\clearpage

\section{Appendix D: Normalized 3K-Parameter Configurations}\label{sec:appendix-d-normalized-3k-parameter-configurations}

{\def\LTcaptype{table} % do not increment counter
\begin{longtable}[]{@{}
  >{\raggedright\arraybackslash}p{(\linewidth - 6\tabcolsep) * \real{0.2500}}
  >{\raggedright\arraybackslash}p{(\linewidth - 6\tabcolsep) * \real{0.2500}}
  >{\raggedright\arraybackslash}p{(\linewidth - 6\tabcolsep) * \real{0.2500}}
  >{\raggedright\arraybackslash}p{(\linewidth - 6\tabcolsep) * \real{0.2500}}@{}}
\toprule\noalign{}
\begin{minipage}[b]{\linewidth}\raggedright
Model
\end{minipage} & \begin{minipage}[b]{\linewidth}\raggedright
Parameters
\end{minipage} & \begin{minipage}[b]{\linewidth}\raggedright
Window
\end{minipage} & \begin{minipage}[b]{\linewidth}\raggedright
Hidden / Architecture
\end{minipage} \\
\midrule\noalign{}
\endhead
\bottomrule\noalign{}
\endlastfoot
MLP-3K & 3,004 & 11 & hidden=231 \\
Hybrid-3K & 3,004 & 11 & hidden=231 (+ DF switch) \\
VAE-3K & 3,037 & 11 & latent=10, hidden=(44, 20) \\
cGAN-3K & 3,004 & 11 & noise=8, g\_hidden=(30, 30, 16), c\_hidden=(32, 16) \\
Transformer-3K & 3,007 & 11 & d\_model=18, heads=2, layers=1 \\
Mamba-S6-3K & 3,027 & 11 & d\_model=16, d\_state=6, layers=1 \\
Mamba-2-3K & 3,004 & 11 & d\_model=15, d\_state=6, chunk\_size=8, layers=1 \\
\end{longtable}
}

All 3K configurations use a window size of 11 (vs.~5 for the original MLP/Hybrid and 7 for the original VAE/cGAN); the original sequence models already used 11.
The parameter counts are within ±1.2\% of the 3,000 target.

\section{Appendix E: Master Experiment Ledger}\label{sec:appendix-e-master-experiment-ledger}

Every numerical table and figure in this thesis traces to exactly one committed result file, verified cell-by-cell against that file by \texttt{verify\_thesis\_tables.py} in the code repository (Section~\ref{sec:reproducibility-and-weight-management}). This ledger lists the modulation, channel, and result source for each major study, grouped in the order they appear. Result-source filenames matching \texttt{e6\_*.npy} are relative to \texttt{e6\_unknown\_channel\_results/}.

{\footnotesize
\begin{longtable}[]{@{}p{0.19\linewidth}p{0.11\linewidth}p{0.13\linewidth}p{0.09\linewidth}p{0.40\linewidth}@{}}
\caption[Master experiment ledger]{Master experiment ledger: modulation, channel, and result-file provenance for every major study.}\label{tbl:ledger}\tabularnewline
\toprule\noalign{}
Study & Modulation & Channel & Tbl/Fig & Result source \\
\midrule\noalign{}
\endfirsthead
\toprule\noalign{}
Study & Modulation & Channel & Tbl/Fig & Result source \\
\midrule\noalign{}
\endhead
\bottomrule\noalign{}
\endlastfoot
Canonical benchmark & QPSK & Rayleigh & tbl:2, fig:10 & \url{results/modulation/qpsk_rayleigh.json} \\
Normalized 3K-param.\ comparison & QPSK & Rayleigh & tbl:8 & \url{results/normalized_3k/3k_rayleigh.json} \\
Coded block-DF, $K$-sweep & QPSK & Rayleigh & tbl:34, 35 & \url{results/coded_df_experiment.json} \\
Coded high-SNR paired re-test & QPSK & Rayleigh & tbl:37 & \url{results/coded_high_budget_test.json} \\
Coded relay-error mechanism & QPSK & Rayleigh & tbl:38 & \url{results/coded_error_mechanism.json} \\
Coded soft-decision relaying & QPSK & Rayleigh & tbl:39 & \url{results/coded_soft_decision.json} \\
Equal-throughput \& latency & QPSK, 16-QAM & Rayleigh & tbl:40, 41 & \url{results/coded_latency_throughput.json} \\
Adaptive modulation/coding (AMC) & QPSK, 16-QAM & Rayleigh & tbl:42 & \url{results/coded_rate_adaptation.json} \\
Capacity under a latency budget & QPSK, 16-QAM & Rayleigh & tbl:43 & \url{results/coded_latency_capacity.json} \\
Unknown ISI (layer 2) & BPSK & AWGN (2nd hop) & tbl:E6 & \url{e6_sim_ported_results.npy} + \url{e6_viterbi_*.npy} \\
Flat-channel control & BPSK & Rayleigh & tbl:E6flat & \url{e6_flat_ported_results.npy} \\
Unknown ISI, QPSK generalization & QPSK & AWGN & tbl:E6qpsk & \url{e6_qpsk_unknown_channel_results.npy} \\
Composite impairment cascade & DBPSK & Rayleigh & E6composite & \url{e6_composite_ported_results.npy} \\
Partial posterior (pilot sweep) & BPSK & Rayleigh & E6partial & \url{e6_partial_ported_results.npy} \\
Blind channel (layer 4) & BPSK & Rayleigh & E6blind & \url{e6_blind_ported_results.npy} \\
Four-layer summary & mixed (above) & mixed (above) & tbl:layers & combination of the above \\
\end{longtable}
}

The result files above are themselves generated by named, committed scripts (e.g.\ \url{coded_rate_adaptation.py} produces \url{coded_rate_adaptation.json}); the correspondence between script and output file is one-to-one and is not separately tabulated here since the filenames make it explicit.

\begin{center}\rule{0.5\linewidth}{0.5pt}\end{center}

%% ── Hebrew Abstract (Back Matter) ───────────────────────────
